\documentclass[
  degree=master,
  degree-type=academic
]{gzuthesis}

\begin{document}

% ==========================================================================
% Front matter
% ==========================================================================

\frontmatter

% --------------------------------------------------------------------------
% 中文摘要
% --------------------------------------------------------------------------

\chapter*{摘\quad 要}
\addcontentsline{toc}{chapter}{摘要}

随着无人机技术、人工智能技术以及低空经济的快速发展，无人机已经广泛应用于
基础设施巡检、应急救援、地质灾害监测、物流运输和环境感知等领域。
无人机飞行过程具有典型的多变量、强耦合和动态时序特征，
其飞行安全受到复杂环境扰动、传感器异常、执行机构故障以及飞控系统状态退化等多种因素影响。
因此，如何充分挖掘无人机飞行数据中的时间依赖关系和变量关联特征，
实现复杂运行环境下无人机飞行状态的准确分析与异常检测，
对于提升无人机系统运行安全性具有重要意义。

本文围绕无人机飞行数据的智能分析与异常检测问题开展研究。
首先，对无人机飞行过程中的多源时序数据进行统一建模，
分析不同飞行变量之间的动态关联特征；
其次，研究基于深度学习的无人机飞行状态表示方法，
通过时间维度和变量维度的联合建模提取飞行数据中的关键动态特征；
进一步，针对不同无人机平台之间数据分布差异较大的问题，
研究具有跨平台迁移能力的飞行状态建模与异常检测方法；
最后，通过公开无人机飞行数据集和实际飞行数据开展实验验证，
分析所提出方法在不同运行条件下的有效性和稳定性。

实验结果表明，所研究的方法能够有效表征无人机正常飞行状态的动态变化规律，
并能够利用飞行状态偏离正常模式的程度识别潜在异常，
为复杂环境下无人机飞行安全监测和智能健康管理提供技术支撑。

\noindent
\textbf{关键词：}
无人机；飞行数据；时间序列；异常检测；深度学习


% --------------------------------------------------------------------------
% English abstract
% --------------------------------------------------------------------------

\chapter*{Abstract}
\addcontentsline{toc}{chapter}{Abstract}

With the rapid development of unmanned aerial vehicles, artificial
intelligence, and the low-altitude economy, UAVs have been increasingly
used in infrastructure inspection, emergency response, geological hazard
monitoring, logistics, and environmental sensing.

UAV flight data are characterized by multivariate coupling, temporal
dependencies, and dynamic operating conditions. The safety of UAV systems
may be affected by environmental disturbances, sensor anomalies,
actuator faults, and degradation of flight-control systems. Therefore,
effective modeling of temporal dynamics and cross-variable dependencies
is important for intelligent UAV condition monitoring.

This thesis investigates intelligent analysis and anomaly detection
methods for UAV flight data. Multisource flight time series are first
represented in a unified form. Deep learning methods are then employed
to characterize temporal evolution and interactions among flight
variables. Furthermore, cross-platform representation learning is
investigated to improve the generalization capability of anomaly
detection models.

\noindent
\textbf{Key words:}
UAV, flight data, time series, anomaly detection, deep learning


% --------------------------------------------------------------------------
% Table of contents
% --------------------------------------------------------------------------

\tableofcontents

% --------------------------------------------------------------------------
% List of figures and tables
% --------------------------------------------------------------------------

\listoffigures
\listoftables


% ==========================================================================
% Main matter
% ==========================================================================

\mainmatter


% ==========================================================================
% Chapter 1
% ==========================================================================

\chapter{绪论}

无人机具有机动灵活、部署快速、运行成本较低以及适应复杂环境能力较强等特点，
近年来已经成为低空智能系统的重要组成部分。
随着无人机应用范围不断扩大，其运行环境和任务类型也日益复杂，
飞行安全问题逐渐成为影响无人机规模化应用的重要因素。

无人机飞控系统需要持续处理惯性测量单元、卫星导航系统、
气压计、磁力计以及执行机构等多个来源的信息。
不同飞行变量之间并非相互独立，而是通过飞行动力学和控制系统形成复杂的动态耦合关系。
当某一传感器、执行机构或控制环节发生异常时，
异常可能首先表现为局部变量偏移，并进一步通过变量关联关系传播到其他飞行状态。


\section{研究背景}

近年来，无人机已经逐渐从传统航拍和娱乐应用向行业应用发展。
在山区基础设施巡检、地质灾害监测和应急救援等场景中，
无人机能够有效降低人工进入危险区域所带来的安全风险。

然而，与规则环境相比，复杂低空环境中的无人机运行状态更加容易受到外界扰动。
例如，山区地形可能引起局部气流变化，
电磁环境可能影响导航和通信系统，
长时间运行也可能导致传感器或执行机构性能下降。

因此，需要利用飞行过程中持续产生的多源时序数据，
实时分析无人机的运行状态，
并及时发现可能影响飞行安全的异常变化。


\subsection{无人机飞行数据特点}

无人机飞行数据通常具有以下特点。

首先，飞行数据具有明显的时间依赖性。
例如，姿态角、角速度和加速度等变量随时间连续变化，
当前飞行状态通常与此前一段时间内的状态密切相关。

其次，不同变量之间存在明显的关联关系。
例如，滚转角变化通常伴随着对应角速度以及控制输入的变化，
执行机构响应也会影响后续飞行状态。

再次，不同无人机平台所记录的变量类型、变量数量和采样频率可能存在差异，
这使得基于单一平台训练的检测模型难以直接应用于其他无人机平台。


\subsection{无人机异常检测需求}

传统无人机故障诊断方法通常针对预定义故障类型构建分类模型。
这种方法需要大量具有明确故障标签的数据。

然而，在实际飞行过程中，
异常数据往往难以大量获取，
某些严重故障甚至不能通过主动飞行试验反复采集。

因此，仅依赖有监督故障分类方法难以覆盖真实环境中可能出现的未知异常。
利用大量正常飞行数据建立正常运行模式，
并根据飞行状态对正常模式的偏离程度进行异常检测，
已经成为具有应用价值的研究方向。


\section{国内外研究现状}

随着深度学习和时间序列分析方法的发展，
基于数据驱动的异常检测方法得到了广泛研究。

\subsection{基于重构的异常检测方法}

基于重构的方法通常利用正常数据训练模型，
使模型学习正常运行状态的数据分布。
在测试过程中，
正常样本通常能够被模型较准确地重构，
而异常样本由于偏离训练数据中的正常模式，
往往产生较大的重构误差。

设输入飞行时间序列为

\begin{equation}
    \mathbf{X}
    =
    \left[
        \mathbf{x}_1,
        \mathbf{x}_2,
        \ldots,
        \mathbf{x}_T
    \right],
\end{equation}

其中，
$T$ 表示时间窗口长度，
$\mathbf{x}_t \in \mathbb{R}^{C}$ 表示第 $t$ 个时刻的
$C$ 维飞行状态。

模型输出重构结果

\begin{equation}
    \hat{\mathbf{X}}
    =
    f_{\theta}(\mathbf{X}),
\end{equation}

其中 $f_{\theta}$ 表示参数为 $\theta$ 的重构模型。

一种常用的异常分数可以表示为

\begin{equation}
    s
    =
    \frac{1}{TC}
    \left\|
        \mathbf{X}
        -
        \hat{\mathbf{X}}
    \right\|_F^2 .
\end{equation}

当异常分数超过给定阈值时，
即可将当前飞行窗口判定为异常状态。


\subsection{基于预测的异常检测方法}

预测类方法利用历史飞行数据预测未来状态。
当系统处于正常状态时，
模型能够根据历史动态模式得到较准确的预测结果；
当系统发生异常后，
实际飞行状态与正常模式下的预测结果之间可能产生明显偏差。

这类方法能够利用飞行数据的时间连续性，
但其检测性能通常受到预测误差累积以及飞行工况变化的影响。


\section{研究内容}

本文围绕无人机飞行数据异常检测问题，
主要开展多变量时间序列建模、正常飞行状态表示、
跨平台泛化以及异常评分等方面的研究。

本文整体研究内容之间的关系如图~\ref{fig:framework} 所示。

\begin{figure}[htbp]
    \centering
    \fbox{
        \parbox[c][4cm][c]{0.75\textwidth}{
            \centering
            此处放置论文总体研究框架图
        }
    }
    \caption{论文总体研究框架}
    \label{fig:framework}
\end{figure}


\section{论文组织结构}

本文共分为若干章节。
第一章介绍研究背景、研究意义以及国内外相关研究现状；
第二章介绍无人机飞行数据和相关理论基础；
第三章研究无人机飞行状态的时序建模方法；
第四章开展相关实验与性能分析；
最后总结全文并讨论后续研究方向。


% ==========================================================================
% Chapter 2
% ==========================================================================

\chapter{无人机飞行数据与相关理论}

本章介绍无人机飞行数据的主要组成、
多变量时间序列基本表示方法以及异常检测问题的基本定义。


\section{无人机飞行数据组成}

无人机飞行过程中通常可以记录姿态、速度、位置、
惯性测量、控制指令和执行机构响应等多种信息。

表~\ref{tab:variables} 给出了部分典型无人机飞行变量。

\begin{table}[htbp]
    \centering
    \caption{典型无人机飞行变量}
    \label{tab:variables}
    \begin{tabular}{lll}
        \toprule
        类别 & 变量 & 说明 \\
        \midrule
        姿态状态 & Roll, Pitch, Yaw & 无人机姿态角 \\
        惯性测量 & Acceleration & 三轴加速度 \\
        惯性测量 & Angular Rate & 三轴角速度 \\
        导航状态 & Position & 无人机位置 \\
        导航状态 & Velocity & 无人机速度 \\
        控制状态 & Control Input & 飞行控制输入 \\
        \bottomrule
    \end{tabular}
\end{table}


\section{多变量时间序列表示}

假设一段无人机飞行数据包含 $C$ 个变量和 $S$ 个采样时刻，
则完整飞行序列可以表示为

\begin{equation}
    \mathbf{X}
    \in
    \mathbb{R}^{S \times C}.
\end{equation}

在实际模型中，
通常通过滑动窗口从完整飞行序列中提取长度为 $T$ 的局部序列，
作为模型输入。


\subsection{数据标准化}

由于不同飞行变量的物理量纲和数值范围存在明显差异，
通常需要在模型训练前进行标准化。

对于变量 $x$，
z-score 标准化可以表示为

\begin{equation}
    x'
    =
    \frac{x-\mu}{\sigma},
\end{equation}

其中 $\mu$ 和 $\sigma$ 分别表示训练数据中的均值和标准差。


\subsection{滑动窗口}

对于完整飞行序列，
采用长度为 $T$、步长为 $d$ 的滑动窗口构建模型输入样本。
该方式能够保持局部飞行动态信息，
并能够将不同长度的飞行记录转换为固定长度输入。


% ==========================================================================
% Chapter 3
% ==========================================================================

\chapter{无人机飞行状态建模方法}

本章研究基于多变量时间序列的无人机飞行状态建模方法，
重点分析时间依赖关系和跨变量关联关系。


\section{总体框架}

模型首先对原始飞行时间序列进行分段和特征映射，
随后利用时序建模模块提取飞行变量在时间维度上的动态特征，
并进一步学习不同变量之间的关联关系。

最终通过重构或预测任务建立正常飞行状态表示，
并根据模型输出误差构造异常分数。


\section{时间特征建模}

设第 $c$ 个飞行变量的窗口序列表示为

\begin{equation}
    \mathbf{x}^{(c)}
    =
    \left[
        x_1^{(c)},
        x_2^{(c)},
        \ldots,
        x_T^{(c)}
    \right].
\end{equation}

时间建模模块用于提取同一变量在不同时间位置之间的相关关系，
从而刻画无人机飞行状态随时间的动态演化过程。


\section{变量关联建模}

除时间依赖关系以外，
无人机不同飞行变量之间还存在明显的物理和控制关联。

例如，当无人机滚转状态发生变化时，
对应角速度、横向加速度以及控制输入通常同时发生变化。
因此，对不同变量之间的动态依赖关系进行建模，
有助于提高模型对异常状态的敏感性。


% ==========================================================================
% Chapter 4
% ==========================================================================

\chapter{实验与结果分析}

本章通过无人机飞行数据实验验证所研究方法的有效性，
并从异常检测性能、模型效率以及不同模块贡献等方面进行分析。


\section{实验设置}

实验数据由正常飞行数据和包含异常状态的测试飞行数据组成。
模型仅利用正常飞行数据进行训练，
在测试阶段根据异常分数判断飞行状态是否出现异常。


\subsection{评价指标}

实验采用准确率、精确率、召回率和 F1 值评价检测性能。

精确率定义为

\begin{equation}
    \mathrm{Precision}
    =
    \frac{TP}{TP+FP},
\end{equation}

召回率定义为

\begin{equation}
    \mathrm{Recall}
    =
    \frac{TP}{TP+FN},
\end{equation}

F1 值定义为

\begin{equation}
    F_1
    =
    \frac{
        2 \times \mathrm{Precision} \times \mathrm{Recall}
    }{
        \mathrm{Precision} + \mathrm{Recall}
    }.
\end{equation}


\section{实验结果}

表~\ref{tab:result} 给出了示例实验结果。

\begin{table}[htbp]
    \centering
    \caption{异常检测实验结果示例}
    \label{tab:result}
    \begin{tabular}{lcccc}
        \toprule
        方法 & Accuracy & Precision & Recall & F1 \\
        \midrule
        Method A & 0.912 & 0.884 & 0.901 & 0.892 \\
        Method B & 0.936 & 0.913 & 0.927 & 0.920 \\
        Proposed & 0.958 & 0.947 & 0.951 & 0.949 \\
        \bottomrule
    \end{tabular}
\end{table}

从示例结果可以看出，
所研究方法能够较好地识别测试序列中的异常状态。


% ==========================================================================
% Chapter 5
% ==========================================================================

\chapter{结论与展望}

本文围绕无人机飞行数据异常检测问题开展研究，
从多变量时间序列建模、正常状态表示和异常评分等方面进行了分析。

研究结果表明，
充分利用无人机飞行变量的时间依赖关系和跨变量关联关系，
有助于提高模型对复杂飞行状态的表征能力。

未来可进一步研究复杂环境因素与无人机飞行状态之间的耦合关系，
增强模型对不同无人机平台和不同运行条件的适应能力，
并探索面向实际飞行任务的在线异常检测和安全决策方法。


% ==========================================================================
% Back matter
% ==========================================================================

\backmatter


% --------------------------------------------------------------------------
% Acknowledgements
% --------------------------------------------------------------------------

\chapter*{致谢}
\addcontentsline{toc}{chapter}{致谢}

在论文研究和撰写过程中，
得到了导师、课题组成员以及相关老师和同学的帮助与支持，
在此表示诚挚感谢。


% --------------------------------------------------------------------------
% References placeholder
% --------------------------------------------------------------------------

\chapter*{参考文献}
\addcontentsline{toc}{chapter}{参考文献}

此处暂作为参考文献排版占位。
后续将采用 BibTeX 或 BibLaTeX 统一管理参考文献，
并按照贵州大学学位论文格式要求进行设置。


% --------------------------------------------------------------------------
% Appendix
% --------------------------------------------------------------------------

\appendix

\chapter{补充实验结果}

本附录给出正文中未展开的部分实验结果和补充说明。


\section{补充实验设置}

这里用于测试附录章节及其编号格式。


\end{document}